%% kapa: exam=1 round=36 year=2025 period=2 form=A subject=appraisal_law subject_ko=감정평가관계법규 seq=31 number=31
%% source: Q-Net 감정평가사 기출문제 (cst00309) artlSeq=5231525 file=raw/1008/5231525/2025년 제36회 감정평가사 1차시험 2교시 문제지 원본.pdf
%% conversion: kordoc 4.24.51 -> tools/gen1.py (auto); review=proofread-sonnet; flags=none
\begin{question}{31}
공간정보의 구축 및 관리 등에 관한 법령상 축척변경에 관한 내용이다. ( )에 들어갈 내용으로 옳은 것은?
\begin{framedtext}
○ 지적소관청은 축척변경을 하려면 축척변경 시행지역의 ( ㄱ ) ( ㄴ ) 이상의 동의를 받아 축척변경위원회의 의결을 거친 후 시ㆍ도지사 또는 대도시 시장의 승인을 받아야 한다. \par
○ 축척변경 시행지역의 토지소유자 또는 점유자는 시행공고가 된 날부터 ( ㄷ )일 이내에 시행공고일 현재 점유하고 있는 경계에 국토교통부령으로 정하는 경계점표지를 설치하여야 한다.
\end{framedtext}
\begin{choices}
  \item ㄱ: 주민, ㄴ: 2분의 1, ㄷ: 20
  \item ㄱ: 주민, ㄴ: 3분의 2, ㄷ: 30
  \item ㄱ: 토지소유자, ㄴ: 2분의 1, ㄷ: 20
  \item ㄱ: 토지소유자, ㄴ: 3분의 2, ㄷ: 30
  \item ㄱ: 토지소유자 또는 점유자, ㄴ: 3분의 2, ㄷ: 20
\end{choices}
\end{question}
\answer{4}
